\documentclass[11pt,a4paper]{article}
\usepackage[utf8]{inputenc}
\usepackage{amsmath,amssymb,amsfonts,amsthm}
\usepackage{geometry}
\usepackage{hyperref}
\usepackage{microtype}

\geometry{margin=1in}

\title{\textbf{Relativistic Future Lookahead \& Global Causality Horizon Engine}\\[0.5em]
\large Mathematical Formalism of Temporal Offsets and Smoothed Time Dilation Transitions}
\author{\textbf{Time Dilation DAW Architecture Team}}
\date{\today}

\begin{document}

\maketitle

\begin{abstract}
In real-time dynamic digital signal processing (DSP), non-causal operations requiring future signal information $\mathbf{x}(t + \Delta t_{\text{future}})$ present a physical impossibility under standard real-time constraints. This paper presents a novel relativistic solution: the \textbf{Global Causality Horizon Offset Engine}. By introducing a dynamic global causality delay $T_{\text{horizon}} = \max_i (\Delta t_{\text{future}, i})$ across the entire workspace, the observer present is shifted backward, allowing future-requesting nodes to inspect the live un-delayed stream. Furthermore, to eliminate audio artifacts and transient phase discontinuities during time dilation factor $\gamma(t)$ modulation, we formulate a 1-pole exponential lowpass filter with user-adjustable cutoff frequency $f_c$ and complete bypass control.
\end{abstract}

\section{Introduction}
Standard digital audio workstations (DAWs) rely strictly on causal signal processing:
\begin{equation}
y(t) = \mathcal{H}\Big( \{x(\tau) \mid \tau \le t\} \Big)
\end{equation}
In relativistic audio synthesis, nodes may request information from future coordinate times $t + \Delta t$. To realize this without violating physical time progression, the engine establishes a global reference frame offset.

\section{Global Causality Horizon Transformation}

In our relativistic framework, \textbf{any time object generating positive time progression} ($\gamma > 1.0$, such as \texttt{[time.warp\~{}]}, \texttt{[time.future\~{}]}, or accelerated \texttt{[time.transport]}) inherently creates a positive time advance relative to standard 1.0x baseline time:
\begin{equation}
\Delta t_{\text{future}, k} = \max\Big(0, \frac{\gamma_k - 1.0}{2}\Big)
\end{equation}

Let $\mathcal{N} = \{N_1, N_2, \dots, N_M\}$ be the graph of audio and control nodes. If nodes generate positive time leads $\Delta t_k > 0$, the master causality horizon $T_{\text{horizon}}$ is dynamically evaluated as:
\begin{equation}
T_{\text{horizon}} = \max_{N_k \in \mathcal{N}} (\Delta t_k)
\end{equation}

The coordinate time transformation for standard (causal) nodes $N_j$ is delayed by $T_{\text{horizon}}$:
\begin{equation}
x_j'(t) = x_j(t - T_{\text{horizon}})
\end{equation}

Conversely, positive time generating nodes $N_k$ access the un-delayed live stream:
\begin{equation}
x_k'(t) = x_k(t) = x_k'(t + T_{\text{horizon}})
\end{equation}

Thus, positive time objects seamlessly project audio and control signals into the future horizon without violating causal execution.

\section{Smoothed Time Dilation Lowpass Filter}

Rapid changes in the relativistic time dilation factor $\gamma(t)$ can introduce steep pitch jumps and phase discontinuities. To guarantee smooth temporal transitions while granting total user control, the effective time dilation factor $\gamma_{\text{smooth}}(t)$ is governed by a first-order continuous-time differential equation:

\begin{equation}
\frac{d \gamma_{\text{smooth}}(t)}{dt} = 2\pi f_c \Big( \gamma_{\text{target}}(t) - \gamma_{\text{smooth}}(t) \Big)
\end{equation}

where $f_c \in [0.1, 50.0] \text{ Hz}$ represents the user-configurable smoothing cutoff frequency.

\subsection{Discrete-Time Recurrence Relation}
In discrete sample domain at sample rate $f_s = 1 / T_s$:
\begin{equation}
\alpha = 1 - \exp\left( -2\pi \frac{f_c}{f_s} \right)
\end{equation}

The smoothed time dilation factor is updated per sample $n$:
\begin{equation}
\gamma_{\text{smooth}}[n] = 
\begin{cases}
\gamma_{\text{target}}[n], & \text{if Bypass is ACTIVE or } f_c = 0 \\
(1 - \alpha)\gamma_{\text{smooth}}[n-1] + \alpha \gamma_{\text{target}}[n], & \text{otherwise}
\end{cases}
\end{equation}

\section{Eliminating Glitches During Horizon Expansion via Lowpass Interpolated Delay}

A critical challenge occurs when the required causality horizon $T_{\text{horizon}}$ increases dynamically (e.g., when a positive time node accelerates from $\gamma = 1.0 \to 3.0$). An instantaneous step jump in delay time $\Delta T$ would cause severe audio buffer tearing, click artifacts, and sample starvation.

To resolve this, our lowpass filter of time $f_c$ is directly applied to the causality horizon expansion rate:

\begin{equation}
\frac{d T_{\text{horizon, smooth}}(t)}{dt} = 2\pi f_c \Big( T_{\text{target}}(t) - T_{\text{horizon, smooth}}(t) \Big)
\end{equation}

During the transition phase ($\frac{dT_{\text{horizon, smooth}}}{dt} > 0$), standard nodes read from the global delay buffer using a \textbf{4-point Hermite fractional delay interpolator}:

\begin{equation}
x_{\text{smooth}}'(t) = \text{Hermite}\Big( x(t - T_{\text{horizon, smooth}}(t)) \Big)
\end{equation}

Because the horizon expansion velocity $\frac{dT_{\text{horizon, smooth}}}{dt} \ll 1.0$ is strictly rate-limited by $f_c$, the delay read pointer moves continuously without buffer tearing. This produces a \textbf{smooth, phase-aligned organic pitch bend} during horizon expansion, completely eliminating harsh glitch spikes while achieving the target future horizon.

If the user activates \textbf{Bypass Mode} ($f_c = 0$), the smoothing constraint is removed, allowing instantaneous quantum step jumps for experimental glitch sound design.

\section{Adaptive Target-Matched Cutoff \& Linear Slew Limiting}

Rather than using a static lowpass cutoff $f_c$, the engine dynamically calculates an \textbf{adaptive cutoff frequency} $f_c^{\text{adaptive}}$ that automatically matches the exact time duration required to travel to the requested future horizon:

\begin{equation}
T_{\text{travel}} = \frac{1}{2} \Big| \gamma_{\text{target}} - \gamma_{\text{smooth}} \Big|
\end{equation}

\begin{equation}
f_c^{\text{adaptive}} = \frac{1}{2\pi \cdot T_{\text{travel}}}
\end{equation}

To complement adaptive lowpass filtering, a \textbf{Linear Slew-Rate Limiter} ($S_{\text{max}}$) bounds the maximum acceleration of coordinate time:

\begin{equation}
\Big| \frac{d\gamma_{\text{smooth}}}{dt} \Big| \le S_{\text{max}}
\end{equation}

This dual-stage architecture guarantees:
\begin{enumerate}
\item \textbf{Constant Doppler Pitch Shift}: Linear slew limiting enforces a constant pitch shift ratio ($\text{ratio} = 1 - S_{\text{max}}$), delivering an organic, musical pitch glide rather than exponential decay.
\item \textbf{Exact Target Arrival}: Adaptive cutoff frequency tuning ensures the system reaches the exact target future horizon $T_{\text{future}}$ precisely when the travel phase completes.
\end{enumerate}

\section{The Future Preparation Horizon: Latent Pre-Rendering and Instant Execution}

The fundamental core principle of our future horizon engine can be stated simply:
\begin{quote}
\textit{Sound from the future may occur; we do not play it yet, but if we need to, we are instantly ready to do so.}
\end{quote}

Formally, the engine maintains a \textbf{Latent Future Preparation Horizon} $\mathbf{x}_{\text{future}}(t + \tau)$ in global RAM memory. While standard audio execution continues at observer present $t$, future audio calculation threads execute ahead by $\Delta t_{\text{lookahead}}$.

\begin{equation}
\mathbf{x}_{\text{future}}(t) = \mathcal{H}\Big( \{x(\tau) \mid \tau \le t + \Delta t_{\text{lookahead}}\} \Big)
\end{equation}

When a time warp acceleration or future trigger event occurs at $t_{\text{event}}$, no calculation lag or buffer starvation takes place: the future audio stream is \textbf{already pre-rendered and phase-aligned}, enabling zero-latency instant execution.

\section{Conclusion}
This framework provides total mathematical consistency for future lookahead audio prediction and smooth temporal warping, preserving phase coherence while offering instant quantum bypass when raw step transitions are desired.

\end{document}
